\documentclass[10pt,a4paper]{amsart}
\usepackage[margin=19mm]{geometry}
\usepackage{amsmath,amssymb,mathtools}
\usepackage[scr=rsfs,cal=euler]{mathalfa}
\usepackage{xfrac}
\usepackage[unicode,hidelinks,bookmarks=false]{hyperref}
\newcommand{\ie}{i.e.\ }
\newcommand{\cf}{cf.\ }
\newcommand{\resp}{respectively\ }
\newcommand{\Subsection}{\subsection}
\providecommand{\nobreakdash}{\nobreak\hbox{-}}
\setlength{\emergencystretch}{2em}
\multlinegap0pt
\usepackage{bm}
\usepackage{bbm}
\usepackage{dsfont}
\usepackage{xspace}
\usepackage{cancel}
\usepackage{mathtools}
\usepackage{amsthm}
\makeatletter
\@ifundefined{theorem}{\newtheorem{theorem}{Theorem}}{}
\@ifundefined{assumption}{\newtheorem{assumption}[theorem]{Assumption}}{}
\@ifundefined{proposition}{\newtheorem{proposition}[theorem]{Proposition}}{}
\makeatother
\newcommand{\B}{\mathcal{B}}
\newcommand{\D}{\mathrm{D}}
\newcommand{\R}{\mathbb{R}}
\renewcommand{\H}{\mathcal{H}}
\title[Functional Gradient Descent with Adaptive Representations]{Functional Gradient Descent with Adaptive Representations}
\author{Daniel Csillag, Rodrigo Schuller, Pedro Dall'Antonia, Leonidas Guibas, Luiz Velho, Tiago Novello}
\date{Source: arXiv:2606.16926v1 (CC BY 4.0)}
\begin{document}
\maketitle

\noindent\emph{Source and licence.} Functional Gradient Descent with Adaptive Representations. Version arXiv:2606.16926v1 (CC BY 4.0), distributed under the Creative Commons Attribution 4.0 International licence, as recorded in the arXiv licence metadata for this exact version and shown on the abstract page https://arxiv.org/abs/2606.16926v1. Full rights notice: licenses/arxiv-2606.16926-CC-BY-4.0.txt.\par\smallskip

\noindent\textit{Editorial scope.} Counted unit: the Proposition whose source label is \texttt{prop:gradient-bridging-rkhs} in the article (source0.tex lines 789--795, source label \texttt{prop:gradient-bridging-rkhs}), delivered together with the article's own complete original proof of it (source0.tex lines 797--847, one \texttt{proof} environment of 51 lines). No mathematics is written, repaired or re-derived: statement and proof are the article's own text line for line. Required context retained uncounted: assumpt:h-descends-in-b (lines 248--251); assumpt:grad-compat (lines 252--255). Every definition, display and label of those lines is retained unchanged. Omitted from this package and not redistributed: 1--247, 256--788, 796--796, 848--1485. Nothing inside a retained excerpt is omitted or altered: every retained body is delivered exactly as the article's lines are written, so no omission receipt of any kind accompanies this package. Citations: the retained excerpts carry no citation command at all, so no bibliography entry of the article is transcribed and no reference is redistributed. Reference adaptation: none was needed and none was made; every reference inside the delivered unit resolves inside it, so no unresolved reference or citation is produced. Printed slips kept literal: no printed word, symbol, hypothesis or number of the counted statement or of its proof was altered. Layout and class: the article's own class is \texttt{article} with packages including neurips\_2025, inputenc, fontenc, xcolor, hyperref, url, booktabs, amsfonts, nicefrac, microtype, dsfont, amsmath, amssymb, amsthm; the wrapper is the lane's pinned \texttt{amsart} editorial wrapper at 10pt on A4, which restates the theorem-like environments the retained text needs and adds the article's own macro definitions that the retained text needs (\texttt{\textbackslash B}, \texttt{\textbackslash D}, \texttt{\textbackslash H}, \texttt{\textbackslash R}), with extra wrapper packages (bm, bbm, dsfont, xspace, cancel, mathtools). The delivered numbering is the unit's own. Assets: no retained span contains a figure environment, a graphics-inclusion command, a PDF-page inclusion, a table or a TikZ picture; no image file is copied into this package, the package declares no asset dependency and no image file is redistributed with it. Licence: version arXiv:2606.16926v1 of this article is distributed under the Creative Commons Attribution 4.0 International licence, as recorded in the arXiv licence metadata for this exact version and shown on the abstract page https://arxiv.org/abs/2606.16926v1. A full rights notice is delivered at licenses/arxiv-2606.16926-CC-BY-4.0.txt. Characters: every retained excerpt is pure ASCII, so the delivered engine, XeLaTeX with its default Latin Modern text font, reports no missing character.
\begin{assumption}[$\H$ descends in $\B$]\label{assumpt:h-descends-in-b}
    There exists some constant $\alpha>0$ such that, for all $f \in \B$,
    \[ \D L(f; \nabla L(f)) \geq \alpha \lVert \nabla L(f) \rVert_\B^2. \]
\end{assumption}

\begin{assumption}[Gradient compatibility]\label{assumpt:grad-compat}
    There exists a constant $\beta > 0$ such that, for every $f \in \B$,
    \[ \lVert \D L(f; \cdot) \rVert_{\B^*} \leq \beta \lVert \nabla L(f) \rVert_{\B}. \]
\end{assumption}

\begin{proposition}\label{prop:gradient-bridging-rkhs}
    Let $L_\mathrm{REG}$ be the regression functional over $\mathcal{H}_K$, and assume the empirical kernel matrix $K \in \R^{n \times n}$ over the training data is strictly positive definite with minimum eigenvalue $\lambda_{\min}(K) > 0$. Furthermore, assume the kernel is bounded, such that $\sup_{x, x'} \lvert K(x, x') \rvert \leq \kappa$. Then, evaluating the gradients over the Banach space $\mathcal{B} = L^\infty$, the loss $L_\mathrm{REG}$ satisfies:
    \begin{enumerate}
        \item[(i)] \textbf{Gradient compatibility (Assumption~\ref{assumpt:grad-compat}):} with constant $\beta = n \lambda_{\min}^{-1}(K)$.
        \item[(ii)] \textbf{$\mathcal{H}$ descends in $\mathcal{B}$ (Assumption~\ref{assumpt:h-descends-in-b}):} with constant $\alpha = \frac{\lambda_{\min}(K)}{n \kappa^2}$.
    \end{enumerate}
\end{proposition}

\begin{proof}
    As established previously, let $c_i = \ell'(f(X_i), Y_i)$ for $i=1, \dots, n$, such that the continuous extension of the gradient to $\mathcal{B} = L^\infty$ is $\nabla L_\mathrm{REG}(f)(x) = \frac{1}{n} \sum_{i=1}^n c_i K(X_i, x)$. The directional derivative for any $\vec{f} \in L^\infty$ is $\D L_\mathrm{REG}(f; \vec{f}) = \frac{1}{n} \sum_{i=1}^n c_i \vec{f}(X_i)$.

    For gradient compatibility,
    we must show that $\lVert \D L_\mathrm{REG}(f; \cdot) \rVert_{(L^\infty)^*} \leq \beta \lVert \nabla L_\mathrm{REG}(f) \rVert_{L^\infty}$.
    Recall that the dual norm of the directional derivative over $L^\infty$ is exactly the scaled $L_1$ norm of the coefficients:
    \[
        \lVert \D L_\mathrm{REG}(f; \cdot) \rVert_{(L^\infty)^*} = \frac{1}{n} \lVert c \rVert_1.
    \]
    Let $g_X \in \R^n$ be the vector of gradient evaluations at the training points, so $g_X^{(i)} = \nabla L_\mathrm{REG}(f)(X_i)$. In matrix form, $g_X = \frac{1}{n} K c$, which gives $c = n K^{-1} g_X$. 
    By the definition of the supremum norm, the maximum absolute evaluation of the gradient over the entire domain bounds the maximum absolute evaluation at the training points: $\lVert g_X \rVert_\infty \leq \lVert \nabla L_\mathrm{REG}(f) \rVert_{L^\infty}$.

    It then follows:
    \begin{align*}
        \lVert c \rVert_1 
        \leq \sqrt{n} \lVert c \rVert_2 
        &= \sqrt{n} \lVert n K^{-1} g_X \rVert_2 \\
        &\leq n^{3/2} \lVert K^{-1} \rVert_2 \lVert g_X \rVert_2 \\
        &\leq n^{3/2} \lambda_{\min}^{-1}(K) \left( \sqrt{n} \lVert g_X \rVert_\infty \right) \\
        &= n^2 \lambda_{\min}^{-1}(K) \lVert g_X \rVert_\infty.
    \end{align*}
    Substituting this into the dual norm equation and applying $\lVert g_X \rVert_\infty \leq \lVert \nabla L_\mathrm{REG}(f) \rVert_{L^\infty}$ yields:
    \[
        \lVert \D L_\mathrm{REG}(f; \cdot) \rVert_{(L^\infty)^*} \leq \frac{1}{n} \left( n^2 \lambda_{\min}^{-1}(K) \lVert \nabla L_\mathrm{REG}(f) \rVert_{L^\infty} \right) = n \lambda_{\min}^{-1}(K) \lVert \nabla L_\mathrm{REG}(f) \rVert_{L^\infty}.
    \]
    Thus, gradient compatibility holds with $\beta = n \lambda_{\min}^{-1}(K)$.

    To establish that $\mathcal{H}$ descends in $\mathcal{B}$,
    we must show $\D L_\mathrm{REG}(f; \nabla L_\mathrm{REG}(f)) \geq \alpha \lVert \nabla L_\mathrm{REG}(f) \rVert_{L^\infty}^2$.
    Evaluating the directional derivative in the direction of the gradient yields:
    \[
        \D L_\mathrm{REG}(f; \nabla L_\mathrm{REG}(f)) = \frac{1}{n} \sum_{i=1}^n c_i \nabla L_\mathrm{REG}(f)(X_i) = \frac{1}{n} c^\top \left( \frac{1}{n} K c \right) = \frac{1}{n^2} c^\top K c.
    \]
    We now require an upper bound for $\lVert \nabla L_\mathrm{REG}(f) \rVert_{L^\infty}^2$ in terms of $c^\top K c$. Using the triangle inequality and the kernel bound $\kappa$:
    \[
        \lVert \nabla L_\mathrm{REG}(f) \rVert_{L^\infty} = \sup_{x} \left\lvert \frac{1}{n} \sum_{i=1}^n c_i K(X_i, x) \right\rvert \leq \frac{1}{n} \sum_{i=1}^n \lvert c_i \rvert \sup_{x} \lvert K(X_i, x) \rvert \leq \frac{\kappa}{n} \lVert c \rVert_1.
    \]
    Squaring both sides and applying the norm inequality $\lVert c \rVert_1^2 \leq n \lVert c \rVert_2^2$, we have:
    \[
        \lVert \nabla L_\mathrm{REG}(f) \rVert_{L^\infty}^2 \leq \frac{\kappa^2}{n^2} \lVert c \rVert_1^2 \leq \frac{\kappa^2}{n} \lVert c \rVert_2^2.
    \]
    By the Rayleigh quotient for positive definite matrices, $c^\top K c \geq \lambda_{\min}(K) \lVert c \rVert_2^2$, meaning $\lVert c \rVert_2^2 \leq \lambda_{\min}^{-1}(K) c^\top K c$. Substituting this gives:
    \[
        \lVert \nabla L_\mathrm{REG}(f) \rVert_{L^\infty}^2 \leq \frac{\kappa^2}{n \lambda_{\min}(K)} c^\top K c.
    \]
    Rearranging to isolate the quadratic form $c^\top K c$:
    \[
        \frac{1}{n^2} c^\top K c \geq \frac{\lambda_{\min}(K)}{n \kappa^2} \lVert \nabla L_\mathrm{REG}(f) \rVert_{L^\infty}^2.
    \]
    Since the left side is exactly $\D L_\mathrm{REG}(f; \nabla L_\mathrm{REG}(f))$, the condition is satisfied with $\alpha = \frac{\lambda_{\min}(K)}{n \kappa^2}$.
\end{proof}

\end{document}
